El presente trabajo se ha desarrollado en el seno del Grupo de Neurocomputación Biológica de la Universidad Autónoma de Madrid, enmarcado en sus líneas de investigación sobre circuitos neuronales biohíbridos. Estos sistemas, en los que neuronas biológicas y modelos computacionales interactúan en tiempo real, exigen un ajuste minucioso de las sinapsis artificiales que median la comunicación en función de la dinámica que se quiera conseguir. Hasta ahora, la adaptación de estos modelos (en específico de sinapsis químicas graduales) a las características electrofisiológicas particulares de cada preparación biológica se realizaba de forma empírica y manual. Este es un proceso complejo, debido al acoplamiento funcional de sus parámetros y a la gran cantidad de los mismos; y lento, lo que dificulta su uso en experimentos que no permitan estar demasiado tiempo haciendo este ajuste.

La motivación principal de este proyecto nace precisamente de la necesidad de automatizar y agilizar dicha tarea. El objetivo no es solo plantear una prueba de concepto frente al problema del ajuste sináptico, sino proporcionar al grupo de investigación una herramienta funcional y robusta, directamente integrable en sus plataformas experimentales.

\hfill \begin{flushright}Marco Gómez Hernández\end{flushright}\hfill

